\chapter{Neural networks and automatic differentiation}
\label{ch:neural-networks}

\begin{learningobjectives}
After this chapter, you should be able to:
\begin{itemize}
  \item express a feedforward network as compositions of affine maps and nonlinearities;
  \item apply the chain rule on a computational graph;
  \item distinguish automatic differentiation from optimization and learning; and
  \item build a reproducible mini-batch training and artifact contract.
\end{itemize}
\end{learningobjectives}

\begin{chapterconnection}
Linear and logistic models used one affine score. Neural networks compose many
affine maps with nonlinear transformations. The single unit remains useful
mathematics, but modern practice gains more from understanding composition,
gradients, tensors, optimization state, and evaluation than from constructing a
general-purpose framework from scratch.
\end{chapterconnection}

\section{Nonlinear composition creates representation depth}

A dense hidden layer maps
\[
h=\phi(Wx+b),
\]
where $W$ is a matrix, $b$ a vector, and $\phi$ an elementwise or structured
activation. A two-layer network may be written
\[
f(x)=W_2\phi(W_1x+b_1)+b_2.
\]
Without a nonlinear $\phi$, the composition collapses to one affine map.
Rectified linear units produce piecewise-linear functions; smooth activations
have different gradient and saturation behavior.

Shapes are contracts. For a batch $X\in\mathbb{R}^{n\times p}$ and hidden width
$h$, $W_1\in\mathbb{R}^{p\times h}$ yields an $n\times h$ activation. Explicit
shape ledgers prevent broadcasting bugs that are numerically valid but
semantically wrong.

\section{Backpropagation is organized chain rule}

A computational graph records operations and intermediate values. Reverse-mode
automatic differentiation begins with the derivative of a scalar loss and
propagates vector-Jacobian products backward through the graph. Backpropagation
is the efficient application of this rule to layered models.

For scalar $\widehat y=w_2\phi(w_1x+b_1)+b_2$ and loss
$\ell=(\widehat y-y)^2/2$,
\[
\frac{\partial\ell}{\partial w_2}=(\widehat y-y)\phi(a),\qquad
\frac{\partial\ell}{\partial w_1}=(\widehat y-y)w_2\phi'(a)x,
\]
where $a=w_1x+b_1$. A finite-difference check compares these derivatives with
small parameter perturbations. It is a debugging oracle for a tiny calculation,
not an efficient training method.

\conceptcheck{Automatic differentiation returns a gradient without error. Which
parts of the learning problem could still be wrong?}

\section{Autodiff, optimization, and estimation are different}

Autodiff differentiates the implemented scalar objective. An optimizer uses
those derivatives to update parameters. Stochastic gradient descent estimates a
full gradient from mini-batches; momentum and adaptive methods transform the
update history. The statistical procedure includes data order, augmentation,
regularization, stopping, and model selection.

These layers can fail independently. A correct gradient can optimize the wrong
loss. A decreasing training loss can coexist with worsening validation behavior.
A valid validation result can still fail under distribution shift.

\section{Loss interfaces should consume logits}

Numerically stable classification losses generally accept raw logits and
combine normalization with log-likelihood internally. Applying a sigmoid or
softmax first and then using a logits-based loss changes the computation and can
damage stability. Class ordering and output shape must be explicit.

Regression output activation follows target support. An unconstrained real
output may suit a centered quantity; a positive quantity may call for a
transformation or distributional model. Architecture is not a substitute for
target semantics.

\begin{warningbox}
Training and evaluation modes can differ. Dropout and batch normalization are
canonical examples. Inference must set evaluation mode, disable gradient
tracking where appropriate, and preserve the exact preprocessing and parameter
state selected for release.
\end{warningbox}

\section{Regularization and training control}

Weight decay, dropout, augmentation, early stopping, architecture size, and
normalization impose different restrictions. Early stopping is model selection:
the validation trajectory chooses a checkpoint. Its data cannot simultaneously
serve as untouched final evaluation.

Monitor training and validation loss, relevant task metrics, gradient norms,
nonfinite values, learning rate, throughput, and resource use. Checkpoints
should include parameters, architecture configuration, preprocessing reference,
optimizer state when resumption matters, epoch, and version metadata.

\begin{professionalbox}
Prefer a small transparent model that completes the full evidence and release
path over a large model whose data lineage, checkpoint, evaluation mode, or
serving requirements cannot be reconstructed. Scale only after the experiment
and resource hypothesis are explicit.
\end{professionalbox}

\section{Frameworks are part of the reproducibility boundary}

A modern tensor framework supplies device-aware arrays, autodiff, modules,
optimizers, data loading, and serialization. It does not define the target,
prevent leakage, select a valid loss, or guarantee deterministic kernels.

Record framework and accelerator versions, device type, seeds, deterministic
settings, numerical precision, and known nondeterministic operations. Test a
saved artifact by loading it in a fresh process and comparing predictions on
versioned fixtures within declared tolerances.

\begin{misconceptionbox}
\textbf{Misconception:} implementing a neural network entirely from scratch is
the best proof of understanding. \textbf{Correction:} one manual forward and
backward pass is valuable. Professional understanding also requires correct
framework use, evaluation, serialization, tests, resource control, and honest
interpretation.
\end{misconceptionbox}

\begin{workedexample}{a tiny network with a large contract}
A two-layer network predicts later battery capacity from standardized early-
cycle features. Before training, one example is passed through a manual network
and its gradients are checked by finite differences. The production candidate
then uses a framework, mini-batches, a fixed outer split, and early stopping on
an inner validation period.

The selected checkpoint is loaded into a fresh process. A fixture verifies
feature order, dtype, output shape, finite values, and prediction tolerance. The
model card reports variability across seeds and future batches. If the network
does not beat the regularized linear champion by a practical margin, it is not
promoted.
\end{workedexample}

\section{Exercises}
\begin{enumerate}
  \item Give dimensions for every tensor in a two-layer network with batch size
  $n$, $p$ features, $h$ hidden units, and $k$ outputs.
  \item Derive the scalar backpropagation equations above and design a central-
  difference gradient check.
  \item Explain why early stopping spends validation information.
  \item Specify a fresh-process test for a serialized neural artifact, including
  tolerance and incompatible-input cases.
\end{enumerate}

\begin{chapterreview}
Neural networks gain expressive power from nonlinear composition of affine
maps. Backpropagation is reverse-mode chain rule; autodiff calculates gradients
of the implemented graph; optimization uses them; and evaluation supports a
generalization claim. Shapes, modes, checkpoints, devices, seeds, preprocessing,
and serialization are essential parts of the learned artifact.
\end{chapterreview}

\begin{notebooklink}
Lecture 16 notebook 00 develops computation graphs and nonlinear composition.
Companion laboratories implement one checked manual backward pass and then train
and serialize a small network with automatic differentiation.
\end{notebooklink}
